% ============================================================================
% Fichier  : partie5/P05_C19_ingenierie_sociale.tex
% Titre    : Ingénierie Sociale
% Auteur   : MINKA MI NGUIDJOÏ Thierry Emmanuel
% Version  : 1.0 -- novembre 2026
% ============================================================================
% Sources : MITRE ATT&CK (T1566 Phishing) ; principes d'éthique des tests
% d'hameçonnage. Campagne autorisée de l'Opération LOGONE : cas fictif,
% chiffres cohérents.
% ============================================================================

\chapter{Ingénierie Sociale}
\label{chap:social}
\chaptermark{Ingénierie sociale}

\epigraph{%
    « On ne teste pas les personnes pour les prendre en défaut~;\\
    on teste un dispositif pour mieux les protéger. »%
}{Principe du cours}

\niveauChapitre{Intermédiaire}{Chapitre~\ref{chap:droit}~; cours SEC3266,
chapitre sur le facteur humain.}

\begin{boxobjectifs}
\begin{itemize}
    \item Définir l'ingénierie sociale et ses principaux leviers.
    \item Décrire les formes d'hameçonnage.
    \item Conduire une campagne d'hameçonnage autorisée et éthique.
    \item Mesurer et interpréter les résultats sans stigmatiser.
    \item Proposer des mesures techniques et de sensibilisation.
\end{itemize}
\end{boxobjectifs}

\bigskip

% ============================================================================
\section*{Question directrice : combien de salariés de LiZbiZ cliquent sur un lien piégé, et qu'en conclure~?}
% ============================================================================

Avec l'accord de la direction et l'information du personnel sur le principe
d'un test, NYONG SÉCURITÉ envoie un courriel d'hameçonnage simulé à
60~salariés. Le but n'est pas de piéger des personnes, mais de mesurer un
risque et d'améliorer un dispositif. Les chiffres orientent la
sensibilisation, non les sanctions.

\medskip
\textbf{Plan du chapitre.}
\begin{enumerate}
    \item Les leviers de l'ingénierie sociale (section~\ref{sec:c19-leviers}).
    \item Les formes d'hameçonnage (section~\ref{sec:c19-formes}).
    \item Une campagne éthique (section~\ref{sec:c19-campagne}).
    \item Mesurer sans stigmatiser (section~\ref{sec:c19-mesure}).
    \item Les parades (section~\ref{sec:c19-parades}).
\end{enumerate}

\separateur

% ============================================================================
\section{Les Leviers de l'Ingénierie Sociale}
\label{sec:c19-leviers}
% ============================================================================

\begin{boxdef}
\termecle{Ingénierie sociale}~: manipulation d'une personne pour obtenir une
information, un accès ou une action, en exploitant la confiance, l'autorité,
l'urgence, la peur ou la serviabilité.
\end{boxdef}

\begin{tableaucours}{L{3.0cm}L{9.2cm}}{%
    \tableauHeader{Levier} &
    \tableauHeader{Exemple}
}
Autorité & Un courriel qui se présente comme la direction \\
Urgence & « Votre compte sera fermé dans une heure » \\
Peur & « Activité suspecte détectée, cliquez pour vérifier » \\
Serviabilité & Un faux collègue demande de l'aide \\
Appât & Une pièce jointe « bulletin de paie » \\
\end{tableaucours}
\vspace{-0.8em}
\begin{center}{\small\itshape Tableau~19.1. Leviers de l'ingénierie sociale}\end{center}

% ============================================================================
\section{Les Formes d'Hameçonnage}
\label{sec:c19-formes}
% ============================================================================

\begin{tableaucours}{L{3.6cm}L{8.6cm}}{%
    \tableauHeader{Forme} &
    \tableauHeader{Description}
}
Hameçonnage de masse & Courriel identique envoyé à un grand nombre \\
Hameçonnage ciblé & Courriel personnalisé visant une personne précise \\
Hameçonnage par message court & Par message téléphonique \\
Hameçonnage vocal & Par appel téléphonique \\
\end{tableaucours}
\vspace{-0.8em}
\begin{center}{\small\itshape Tableau~19.2. Formes d'hameçonnage (ATT\&CK T1566)~\cite{mitreattack}}\end{center}

% ============================================================================
\section{Une Campagne Éthique}
\label{sec:c19-campagne}
% ============================================================================

\begin{boxprocedure}
\textbf{Conduire une campagne d'hameçonnage autorisée}
\begin{enumerate}
    \item Accord écrit de la direction~; information préalable du personnel
          sur le principe (pas sur la date).
    \item Scénario réaliste mais sans contenu humiliant ni thème sensible.
    \item Page d'atterrissage pédagogique~: la personne qui clique apprend
          aussitôt qu'il s'agissait d'un test.
    \item Collecte limitée~: taux de clic et de saisie, jamais le mot de
          passe réel.
    \item Résultats agrégés~; aucune liste nominative de « fautifs ».
\end{enumerate}
\end{boxprocedure}

\begin{boxalert}
\textbf{Violation typique.} Diffuser la liste des personnes qui ont cliqué.

\smallskip
\textbf{Règle~:} une campagne mesure un dispositif, pas des individus. On ne
collecte jamais les mots de passe réels saisis~; on ne nomme personne~; on
ne tourne pas le test en sanction, sous peine de détruire la confiance et le
signalement.
\end{boxalert}

% ============================================================================
\section{Mesurer sans Stigmatiser}
\label{sec:c19-mesure}
% ============================================================================

\begin{boxscenario}{Opération LOGONE. La campagne de 60 courriels}
\textbf{Constatations.} 60~courriels envoyés. 14~personnes ont cliqué sur le
lien (23,3~\%). 5~ont commencé à saisir des identifiants sur la page de test
(la saisie a été interrompue, aucun mot de passe réel n'a été conservé).
9~personnes ont signalé le courriel à l'informatique.

\medskip
\textbf{Analyse} (confiance élevée). Le taux de clic (23,3~\%) rejoint le
constat T-11. Le taux de signalement (15~\%) est encourageant et constitue
le meilleur indicateur à faire progresser. Aucune conclusion individuelle
n'est tirée.
\end{boxscenario}

\begin{boxobjet}{T-11 · Taux de clic élevé à l'hameçonnage simulé}
\textbf{Description~:} 14~clics sur 60 (23,3~\%) lors de la campagne
autorisée~; 5~débuts de saisie. \textbf{Preuve~:} statistiques agrégées de
la plateforme de test. \textbf{Gravité~:} moyenne. \textbf{Impact~:} un
courriel piégé réel ouvrirait un premier accès. \textbf{Recommandation~:}
sensibilisation régulière, bouton de signalement, second facteur,
marquage des courriels externes. \textbf{Détection~:} suivi du taux de
signalement et du taux de clic dans le temps.
\end{boxobjet}

% ============================================================================
\section{Les Parades}
\label{sec:c19-parades}
% ============================================================================

\begin{tableaucours}{L{4.4cm}L{7.8cm}}{%
    \tableauHeader{Parade} &
    \tableauHeader{Effet}
}
Sensibilisation régulière & Réduit le taux de clic, augmente le signalement \\
Bouton de signalement & Transforme chaque salarié en capteur \\
Second facteur & Un mot de passe saisi par erreur ne suffit plus \\
Marquage des courriels externes & Aide à repérer l'usurpation \\
Filtrage amont & Bloque une partie des courriels piégés \\
\end{tableaucours}
\vspace{-0.8em}
\begin{center}{\small\itshape Tableau~19.3. Parades à l'hameçonnage}\end{center}

\begin{boxmetier}{RSSI et analyste SOC}
\textbf{Faire du salarié un capteur.} L'objectif d'une campagne n'est pas de
faire baisser le taux de clic à zéro, impossible, mais de faire monter le
taux de signalement~: un salarié qui signale vite permet au SOC de réagir
avant la propagation.
\end{boxmetier}

\separateur

\begin{boxresume}
\textbf{Ce qu'il faut retenir de ce chapitre~:}
\begin{itemize}
    \item \textbf{L'ingénierie sociale} exploite autorité, urgence, peur,
          serviabilité, appât.
    \item \textbf{L'hameçonnage} prend plusieurs formes, de la masse au
          ciblé.
    \item \textbf{Une campagne éthique} informe, ne collecte pas les mots de
          passe, ne nomme personne.
    \item \textbf{On mesure un dispositif}, pas des individus.
    \item \textbf{Le taux de signalement} est le meilleur indicateur à faire
          progresser.
\end{itemize}
\end{boxresume}

\begin{boxinfo}
\textbf{Points de contrôle de l'élève.} Je sais~: définir l'ingénierie
sociale et ses leviers~; citer les formes d'hameçonnage~; conduire une
campagne éthique~; interpréter les résultats sans stigmatiser~; proposer des
parades.
\end{boxinfo}

\bigskip

\begin{boxexo}[title={\ding{45}\quad Exercice~19.1. Leviers \niveaubase}]
Associez à chaque courriel le levier exploité~: « la direction exige »~;
« compte fermé dans une heure ».
\end{boxexo}

\begin{boxexo}[title={\ding{45}\quad Exercice~19.2. Formes \niveaubase}]
Distinguez hameçonnage de masse et hameçonnage ciblé.
\end{boxexo}

\begin{boxexo}[title={\ding{45}\quad Exercice~19.3. Taux \niveaubase}]
Calculez le taux de clic et le taux de signalement de la campagne.
\end{boxexo}

\begin{boxexo}[title={\ding{45}\quad Exercice~19.4. Page pédagogique \niveaumoyen}]
Pourquoi la page d'atterrissage doit-elle être pédagogique et non
culpabilisante~?
\end{boxexo}

\begin{boxexo}[title={\ding{45}\quad Exercice~19.5. Collecte \niveaumoyen}]
Pourquoi ne conserve-t-on jamais les mots de passe réels saisis~?
\end{boxexo}

\begin{boxexo}[title={\ding{45}\quad Exercice~19.6. Signalement \niveaumoyen}]
Pourquoi le taux de signalement est-il un meilleur objectif que le taux de
clic~?
\end{boxexo}

\begin{boxexo}[title={\ding{45}\quad Exercice~19.7. T-11 \niveaumoyen}]
Rédigez le plan d'action lié à T-11.
\end{boxexo}

\begin{boxexo}[title={\ding{45}\quad Exercice~19.8. Thèmes sensibles \niveaumoyen}]
Pourquoi évite-t-on les thèmes sensibles (salaire, santé) dans un scénario
de test~?
\end{boxexo}

\begin{boxexo}[title={\ding{45}\quad Exercice~19.9. Campagne \niveauexpert}]
Rédigez la note d'information préalable du personnel annonçant le principe
d'une campagne, sans en révéler la date.
\end{boxexo}

\begin{boxexo}[title={\ding{45}\quad Exercice~19.10. Programme \niveauexpert}]
Proposez un programme annuel de sensibilisation et de mesure pour LiZbiZ.
\end{boxexo}

\bigskip

\begin{boxinfo}
\textbf{Pour aller plus loin.}
\begin{itemize}
    \item MITRE ATT\&CK~\cite{mitreattack}.
\end{itemize}
\end{boxinfo}

\begin{boxinfo}
\textbf{Transition.} L'ensemble de ces techniques se combine dans un
exercice réaliste~: le Chapitre~\ref{chap:rouge} traite de l'équipe rouge et
de l'équipe violette.
\end{boxinfo}
